\documentclass{beamer}
\usetheme{Madrid}
\usepackage{xcolor}
\usepackage{listings}

\usepackage{tikz}
\usetikzlibrary{calc}
\usepackage{graphicx}

\title{ARIN7102 - Project 3 \\Online Health Science Knowledge Chatbot}
% \subtitle{Methods and Technologies for Sub-Tasks B/C, E/F, L/M}
\author{\textbf{\\Group 3.2}}
\date{\today}

\begin{document}
	
	\frame{\titlepage}
	
	\begin{frame}{Project Introduction}
		\begin{itemize}
			\item \textbf{Core Objective}: Develop an intelligent chatbot for parsing health science queries through multi-source data mining (\textit{Healthline} pages, MedQA datasets, web sources) and various NLP techniques
			\item \textbf{Technical Architecture \& Processing Pipeline}: 
			\begin{itemize}
				\item \textcolor{blue}{Vector base construction} via keyword-knowledge mapping
				\item QA pairs' topic \textcolor{blue}{classification}
				\item Most relevant pages \textcolor{blue}{retrieval}
				\item \textcolor{blue}{LLM+RAG} to summarize the answers
				\item \textcolor{blue}{Predictive modeling} for anticipating follow-up questions
			\end{itemize}
			
			\item \textbf{Expected Outcomes}:
				\begin{itemize}
					\item Real-time symptom analysis with differential diagnosis
					\item Adaptive response generation through ensemble modeling
					\item Encapsulate all functions into an interactive visual chatbot
				\end{itemize}
		\end{itemize}
	\end{frame}
	
	\begin{frame}{Content}
	\begin{itemize}
		\item \textbf{1. Data Processing}
		\begin{itemize}
			\item 1.1 Preprocessing
			\item 1.2 Vectorization
			\item 1.3 Topic Clustering
		\end{itemize}
		
		\item \textbf{2. Answer Generation}
		\begin{itemize}
			\item 2.1 RAG Implementation
			\item 2.2 Hybrid Search Components
		\end{itemize}
		
		\item \textbf{3. Next Question Prediction}
		
		\item \textbf{4. Dashboard/WebApp Development}
		
		\vspace{0.09234cm}
		\item \large{*Health Science Chatbot Demenstration Video}
	\end{itemize}
	\end{frame}
	
	% 数据预处理模块
	\begin{frame}{1.1 - Preprocessing}
		\footnotesize
		\textbf{Implemented Components}
		\begin{itemize}
			\item \textbf{Input Sources}:
			\begin{itemize}
				\item Healthline articles (CSV: title/content) (Fetched from \underline{\textit{healthline.com}})
				\item MedQA datasets (CSV: question/answer) (From \textit{HuggingFace})
			\end{itemize}
			
			\item \textbf{Processing Flow}:
			\begin{itemize}
				\item Text cleaning: HTML removal, punctuation normalization
				\item Keyword Extraction:
				\begin{itemize}
					\item             Use NLTK for tokenization
					\item             Remove stopwords (including general and medical domain-specific stopwords)
					\item Use TF-IDF algorithm to extract important keywords
				\end{itemize}
				
				\item QA pair generation: Title→Question mapping
			\end{itemize}
			
			\item \textbf{Output}:
			\begin{itemize}
				\item Standardized QA pairs 
				\item keyword index
			\end{itemize}
		\end{itemize}
	\end{frame}
	
	\begin{frame}[t]{1.2 - Vectorization}
		\begin{itemize}
			\item Use \textcolor{blue}{BERT} to vectorize (Vector dimension: 384)
			\item Use a \textcolor{blue}{sliding window} to dynamically block long text (Window length: 512; Overlap rate: 0.2)
			\item Use \textcolor{blue}{ChromaDB} to store data: Keywords are stored by String to preserve more effective information
		\end{itemize}
		
		\tikz[remember picture,overlay] \node[anchor=north east] at ($(current page.north east)-(1.77cm,4.17cm)$) {\includegraphics[width=9cm]{syc\_1}};
	\end{frame}
	
	% 主题聚类模块
	\begin{frame}{1.3 - Topic Clustering}
		\footnotesize
		\textbf{Technical Implementation}
		\begin{itemize}
			\item \textbf{Dimensionality Reduction}:
			\begin{itemize}
				\item UMAP: n\_components=50, cosine metric
				\item Variance retained: 89.7\%
			\end{itemize}
			
			\item \textbf{Clustering Algorithm}:
			\begin{itemize}
				\item HDBSCAN: min\_cluster\_size=100
				\item Cluster selection: EOM method
			\end{itemize}
			
			\item \textbf{Feature}:
			\begin{itemize}
				\item Train random forest classifier based on clustering results
				\item Automatic integration with vector database system
			\end{itemize}
		\end{itemize}
	\end{frame}
	
		\begin{frame}{2.1 - RAG Implementation}
		\footnotesize
		\begin{enumerate}
			\item \textbf{Hybrid Retrieval}
			\begin{itemize}
				\item Boolean Search: Keywords + Inverted Index (TF-IDF)
				\item Semantic Search: BioBERT Embeddings + Cosine Similarity
				\item Fusion: Combine Results (40\% Boolean + 60\% Semantic)
			\end{itemize}
			
			\item \textbf{Answer Generation}
			\begin{itemize}
				\item Template: Symptoms → Causes → Treatments
				\
			\end{itemize}
			
			\item \textbf{Quality Control}
			\begin{itemize}
				\item Deduplication: Merge Similar Results
				\item Validation: Check Medical Accuracy
			\end{itemize}
		\end{enumerate}
	\end{frame}
	
	\begin{frame}{2.2 - Hybrid Search Components}
		\centering
		\includegraphics[width=0.8\textwidth]{hybrid\_search\_flow}
		\begin{itemize}
			\item \textcolor{blue}{Step 1:} Boolean Filtering (Fast)
			\item \textcolor{blue}{Step 2:} Semantic Ranking (Accurate)
			\item \textcolor{blue}{Step 3:} Weighted Fusion
		\end{itemize}
	\end{frame}
	
	% Most Relevant Pages Retrieval
	\begin{frame}{2.2 - Hybrid Search Components (Cont.)}
		\footnotesize
		\textbf{Objective}: Designed a hybrid retrieval system combining \textcolor{blue}{keyword matching} with query's \textcolor{blue}{semantic analysis} for precision-focused health knowledge pages retrieval
		
		\textbf{Technical Stack}:
		\begin{itemize}
			\item \textbf{Libraries}: NLTK, Sentence Transformers, spaCy (SciSpaCy), etc.
			\item \textbf{Models}: MPNet (Sentence-BERT), BM25, and TF-IDF
		\end{itemize}
		
		\textbf{Pipeline}:
		\begin{enumerate}
			\item \textcolor{blue}{Boolean filtering} via keyword inverted index
			\item \textcolor{blue}{Semantic scoring} using MPNet embeddings
			\item \textcolor{blue}{Weighted fusion} of TF-IDF, BM25, and medical\\ entity scores
		\end{enumerate}
		
		\textbf{Innovations}:
		\begin{itemize}
			\item \textbf{Context-Aware Weighting}: Combined entity recognition (spaCy patterns) with TF-IDF statistical relevance and BM25 term frequency
			\item \textbf{Hybrid Retrieval}: Achieved 82\% precision by merging logical matching results with output semantic ranking (Top 10 relevance)
			\item \textbf{Dynamic Thresholding}: Adaptive weight allocation for clinical terms (e.g., "migraine triggers" 2.0x weight vs general terms)
		\end{itemize}
		
		\tikz[remember picture,overlay] \node[anchor=north east] at ($(current page.north east)-(0.12cm,3.17cm)$) {\includegraphics[trim=0cm 3cm 0cm 7cm, clip, width=4.9cm]{Hybrid Retrieval Pipeline Flowchart}};
	\end{frame}
	
	\begin{frame}{3 - Next Question Prediction}
		\footnotesize
		\begin{enumerate}
			\item \textbf{Input Formatting}
			\begin{itemize}
				\item Context: [Current Question + Answer]
				\item Prompt: "Suggest 3 logical follow-up questions about:"
			\end{itemize}
			
			\item \textbf{LLM Configuration}
			\begin{itemize}
				\item Pre-trained Model: QWEN-2.5-7B-chat
				\item Generation Parameters:
				\begin{itemize}
					\item Temperature: 0.7 (creativity control)
					\item Max Length: 50 tokens
					\item Beam Search: Keep top 5 candidates
				\end{itemize}
			\end{itemize}
			
			\item \textbf{Output Processing}
			\begin{itemize}
				\item Filter medical terms via UMLS dictionary
				\item Remove duplicate/redundant questions
			\end{itemize}
		\end{enumerate}
	\end{frame}
	
	\begin{frame}[t]{4 - Dashboard/WebApp}
		\begin{itemize}
			\item Use native \textcolor{blue}{HTML/JS/CSS} to complete the frontend
			\item Use \textcolor{blue}{SQLite3} to store user information and historical sessions
			\item Use \textcolor{blue}{Flask} to complete backend development, and provide \textcolor{blue}{interfaces} for deploying model
		\end{itemize}
		
		\vspace{3.97cm}
		- Detailed effects could be referred to our following demonstration video.
		
		\tikz[remember picture,overlay] \node[anchor=north east] at ($(current page.north east)-(3.29cm,4.17cm)$) {\includegraphics[width=5.67cm]{syc\_2}};
	\end{frame}
	
	\begin{frame}
		\begin{center}
			\LARGE{\textbf{Thank you!}}
		\end{center}
	\end{frame}
	
\end{document}